\scanpage{76}
Si $S$ est de type fini sur $\mathbb{F}_p$, on % 1
aura (travaillant avec des \uline{faisceaux % 2
étales}) si $f:S\longrightarrow\operatorname{Spec}\mathbb{F}_p$ % 3
\[
\begin{aligned}
0\to H^1(\mathbb{F}_p,R^0f_*A)\to H^1(S,A)\longrightarrow {}&
 H^0(\mathbb{F}_p,\mathbb{R}^1f_*(A))\longrightarrow0
 \deleted{\unsure{H^2(\mathbb{F}_p,R}}\\ % 4
&\rotatebox{90}{$=$}\\
&H^1(\bar S,\bar A)^\pi % 5
\end{aligned}
\]
Or on a, \unsure{tjrs|tjous|toujours} par \textquotedblleft Kummer\textquotedblright % 6
\[
\begin{aligned}
&\quad\inserted{M\otimes\mathbb{Q}_\ell/\mathbb{Z}_\ell}\\ % 7
0\to{}&H^0(\bar S,\bar A)\unsure{\otimes}\mathbb{Q}_\ell/\mathbb{Z}_\ell
 \;\unsure{-|\to}\; H^1(\bar S,{}_{\ell^\infty}\bar A)
 \to H^1(\bar S,\bar A)(\ell)\to0 % 8
\end{aligned}
\]
d'où le fait que $H^1(\bar S,\bar A)(\ell)$ est de % 9
poids $1-1=0$ \inserted{si p.ex. $S$ est projectif lisse sur $\mathbb{F}_p$} -- à priori rien ne s'oppose % 10
à ce qu'il y ait un gros module % 11
d'invariants. D'autre part, on a % 12
\[
\begin{aligned}
0\to &A^{\text{\unsure{fixe|fix}}}\longrightarrow R^0f_*(A)\longrightarrow M\longrightarrow 0\\ % 13
&\text{partie fixe}\hspace{29mm}\pi\text{-module,}\\ % 14
&\text{de la V.A. }A\text{ sur }S\hspace{12mm}\text{de type fini sur }\mathbb{Z}
\end{aligned} % 15
\]
d'où % 16
\[
\begin{aligned}
H^1(\mathbb{F}_p,A^{\text{fixe}})&\longrightarrow H^1(\mathbb{F}_p,\mathbb{R}^0f_*\unsure{M|A})
 \longrightarrow H^1(\mathbb{F}_p,M)\to0\\ % 17
\rotatebox{90}{$=$}\qquad&\hspace{46mm}\rotatebox{90}{$=$}\\
0\qquad&\hspace{35mm}(\unsure{\underline M|M}\otimes\mathbb{Q}/\mathbb{Z})^\pi\\ % 18
\text{par Lang}\qquad&
\end{aligned} % 19
\]
% archive sheet number: 75
